\lecture{17}
\title{$p$-Values}

\begin{document}

Note that the behavior of a test $\Psi: \{\mathrm{data}\} \to \{0, 1\}$
is a function of the level $\alpha$ that is set.

\textbf{Definition.}
Given a test $\Psi$ and some inputted data $\{X_i\}_{i = 1}^n$,
the \textbf{$p$-value} of $\Psi$ is the smallest level $\alpha$
at which $\Psi$ rejects $H_0$.

Suppose we perform a test and get a $p$-value of\dots
\begin{itemize}
    \item \dots less than $1\%$.
    Then we say we have \emph{very strong} evidence against $H_0$.
    \item \dots between $1\%$ and $5\%$.
    Then we say we have \emph{strong} evidence against $H_0$.
    \item \dots between $5\%$ and $10\%$.
    Then we say we have \emph{weak} evidence against $H_0$.
    \item \dots greater than $10\%$.
    Then we say we have \emph{little or no} evidence against $H_0$.
\end{itemize}

\begin{remark}
    In most settings,
    we can simplify the definition of a $p$-value
    to just be ``the probability of observing results as
    or more extreme than what we did,
    assuming that $H_0$ was true''.
\end{remark}

\textbf{Example.} Suppose we observe $\bar{X}_n = 33.4$ with $n = 164$ and $\hat{\sigma} = 12$,
and we are doing a test on $H_0: \mu \leq 30$ and $H_1: \mu > 30$.
Using the Wald test $\Psi$, our test statistic is:
\[ T_n = \frac{\sqrt{n}}{\hat{\sigma}}(\bar{X}_n - 30)
= \frac{\sqrt{164}}{12}(33.4 - 30) \approx 3.63. \]
The test $\Psi$ rejects $H_0$ if $3.63 \geq z_{\alpha}$.
The least $\alpha$ for which this occurs is $1 - \Phi(3.63) \approx 0.01\%$,
so the $p$-value is $0.01\%$, and we have very strong evidence against $H_0$.

\end{document}